% ------------------------------------------------------------------------------
% Este fichero es parte de la plantilla LaTeX para la realización de Proyectos
% Final de Grado, protegido bajo los términos de la licencia GFDL.
% Para más información, la licencia completa viene incluida en el
% fichero fdl-1.3.tex

% Copyright (C) 2012 SPI-FM. Universidad de Cádiz
% ------------------------------------------------------------------------------

\documentclass[a4paper,11pt]{book}

% PAQUETES
\usepackage{./estilo/paquetes}
\usepackage{./estilo/colores}
\usepackage{./estilo/comandos}

% Ruta al directorio de imágenes
\graphicspath{{./img/}} 

% METADATOS
\title{PLANTILLA PARA TRABAJO FIN DE GRADO}
\author{Profesor Profesor Profesor}
\date{Junio 2012} 
 
\begin{document}

\pagestyle{empty}

% PORTADAS
% ------------------------------------------------------------------------------
% Este fichero es parte de la plantilla LaTeX para la realización de Proyectos
% Final de Grado, protegido bajo los términos de la licencia GFDL.
% Para más información, la licencia completa viene incluida en el
% fichero fdl-1.3.tex

% Copyright (C) 2012 SPI-FM. Universidad de Cádiz
% ------------------------------------------------------------------------------


\begin{titlepage}

  \begin{center}

    \includegraphics[width=0.3\textwidth]{logo-uca.png} \\
    
    \vspace{2.5cm}
    
    \LARGE{\textbf{ESCUELA SUPERIOR DE INGENIERÍA}} \\
    
    \vspace{1.0cm}
    
    \Large{\textbf{GRADO EN INGENIERÍA INFORMÁTICA}} \\
    
    \vspace{3.0cm}
    
    \Large{PLANTILLA PARA PROYECTO DE FIN DE GRADO} \\
    
    \vspace{2.5cm}
    
    \Large{Alumno Alumno Alumno} \\
  
    \vspace{0.5cm}

    \large{\today}
    
  \end{center}
\end{titlepage}

\cleardoublepage

% ------------------------------------------------------------------------------
% Este fichero es parte de la plantilla LaTeX para la realización de Proyectos
% Final de Grado, protegido bajo los términos de la licencia GFDL.
% Para más información, la licencia completa viene incluida en el
% fichero fdl-1.3.tex

% Copyright (C) 2012 SPI-FM. Universidad de Cádiz
% ------------------------------------------------------------------------------


\begin{center}

  \includegraphics[width=0.3\textwidth]{logo-uca.png} \\

  \vspace{2.5cm}

  \Large{ESCUELA SUPERIOR DE INGENIERÍA} \\

  \vspace{1.0cm}

  \large{GRADO EN INGENIERÍA EN INFORMÁTICA} \\

  \vspace{2.0cm}

  \large{PLANTILLA PARA PROYECTO DE FIN DE GRADO} \\

  \vspace{2.5cm}

\end{center}

\begin{itemize}
\item \large{Departamento: Lenguajes y Sistemas informáticos}
\item \large{Director del proyecto: Profesor Profesor Profesor}
\item \large{Autor del proyecto: Alumno Alumno Alumno}
\end{itemize}

\vspace{0.2cm}

\begin{flushright}
  \large{Cádiz, \today} \\

  \vspace{2.5cm}

  \large{Fdo: Alumno Alumno Alumno}
\end{flushright}

\cleardoublepage

% PRELIMINARES
% ------------------------------------------------------------------------------
% Este fichero es parte de la plantilla LaTeX para la realización de Proyectos
% Final de Grado, protegido bajo los términos de la licencia GFDL.
% Para más información, la licencia completa viene incluida en el
% fichero fdl-1.3.tex

% Copyright (C) 2012 SPI-FM. Universidad de Cádiz
% ------------------------------------------------------------------------------

\thispagestyle{empty}

\noindent \textbf{\begin{Large}\textit{Agradecimientos}\end{Large}} 
\newline
\newline
\noindent\textit{Introduzca aquí, si lo desea, los agradecimientos.}

\newpage

% ------------------------------------------------------------------------------
% Este fichero es parte de la plantilla LaTeX para la realización de Proyectos
% Final de Grado, protegido bajo los términos de la licencia GFDL.
% Para más información, la licencia completa viene incluida en el
% fichero fdl-1.3.tex

% Copyright (C) 2012 SPI-FM. Universidad de Cádiz
% ------------------------------------------------------------------------------

\thispagestyle{empty}

\noindent \textbf{\begin{Large}Resumen\end{Large}} 
\newline
\newline
\noindent Introduzca aquí un resumen no superior a 500 palabras, que servirá de descripción pública del trabajo realizado. 
\newline

\noindent {\bf Palabras clave:} Lista de palabras clave que reflejen el contenido del trabajo en aras de facilitar su búsqueda en sistemas bibliográficos.

\newpage


\frontmatter

% INDICES
\tableofcontents
\listoffigures
\listoftables

\mainmatter

% PROLEGÓMENO
\part{Prolegómeno}
\null\vfill
\noindent La primera parte de la memoria del PFC debe contener una introducción y una planificación del proyecto.\\

La introducción es un capítulo que, a modo de resumen, debe contener una breve descripción del contexto de la disciplina en la que el proyecto tiene aplicación y la motivación para su desarrollo, así como del alcance previsto.\\

El segundo capítulo debe incluir una planificación del proyecto. La planificación deberá ajustarse a las prácticas de ingeniería en general, y de la ingeniería del software en particular. Deberá tener en cuenta los plazos, los entregables (documentos y software), los recursos (humanos y de equipamiento inventariable) y el método de ingeniería de software a emplear.
\\

\chapter{Introducción}
% ------------------------------------------------------------------------------
% Este fichero es parte de la plantilla LaTeX para la realización de Proyectos
% Final de Grado, protegido bajo los términos de la licencia GFDL.
% Para más información, la licencia completa viene incluida en el
% fichero fdl-1.3.tex

% Copyright (C) 2012 SPI-FM. Universidad de Cádiz
% ------------------------------------------------------------------------------

A continuación, se describe la motivación del presente proyecto y su alcance. También se incluye un glosario de términos y la organización del resto de la presente documentación.

\section{Motivación}
Qué motivación nos ha llevado a su desarrollo. Contexto y ámbito en el que se desarrolla el proyecto.

\section{Alcance} 
Esta sección debe describir a qué elementos organizativos de la organización Cliente afecta el desarrollo del nuevo sistema. También debe describir los principales objetivos que se esperan alcanzar cuando el sistema a desarrollar esté en producción.

\section{Glosario de Términos} 
Esta sección debe contener una lista ordenada alfabéticamente de los principales términos, acrónimos y abreviaturas específicos del dominio del problema, especialmente de los que se considere que su significado deba ser aclarado. Cada término, acrónimo o abreviatura deberá acompañarse de su definición.

\section{Organización del documento}
Descripción de los contenidos de la presente memoria, así como del software entregado en soporte informático.






\chapter{Planificación}
% ------------------------------------------------------------------------------
% Este fichero es parte de la plantilla LaTeX para la realización de Proyectos
% Final de Grado, protegido bajo los términos de la licencia GFDL.
% Para más información, la licencia completa viene incluida en el
% fichero fdl-1.3.tex

% Copyright (C) 2012 SPI-FM. Universidad de Cádiz
% ------------------------------------------------------------------------------

En esta sección se describen todos los aspectos relativos a la gestión del proyecto: metodología, organización, costes, planificación, riesgos y aseguramiento de la calidad.

\section{Metodología de desarrollo}
Definición del proceso de desarrollo, ciclo de vida y metodología empleada durante la elaboración del proyecto. Las fases y/o iteraciones que proponga el método empleado deberán quedar recogidas en la planificación que se detalle más adelante.

\section{Planificación del proyecto}
Estimación temporal y definición del calendario básico (hitos principales e iteraciones). Desarrollo de la planificación detallada, utilizando un diagrama de Gantt. Los diagramas de Gantt que se vean correctamente (girados y divididos si hace falta).\\

Se debe incluir una comparación cuantitativa del tiempo y el esfuerzo realmente invertido frente al estimado y planificado. Estos datos pueden recogerse del sistema de gestión de tareas empleado para el seguimiento del proyecto.

\section{Organización}
Relación de las personas (roles) involucradas en el proyecto y de cómo se estructuran las relaciones entre las mismas para ejecutar el proyecto. Relación de los recursos inventariables utilizados en el proyecto: equipamiento informático (hardware y software), herramientas empleadas, etc. \\

\section{Costes}
Estudio y presupuesto de los costes de los recursos (humanos y materiales) descritos anteriormente, necesarios para el proyecto.

Para el cálculo de costes de personal pueden consultarse las tablas salariales de la UCA para el personal técnico de apoyo contratado laboral \cite{paslaboral}, o bien otras más ajustadas a la realidad. El cálculo del coste del personal del proyecto debe hacerse en personas-mes, y luego hacer la correspondencia al coste monetario.\\

\section{Riesgos}
Enumeración de los riesgos del proyecto, indicando su posible impacto (efecto que la ocurrencia del citado riesgo tendría en el desarrollo del proyecto) y la probabilidad de ocurrencia. Una vez los riesgos son identificados y priorizados, hay que definir los planes necesarios para reducir los efectos del riesgo una vez se haya materializado o disminuir que este ocurra.\\

\section{Aseguramiento de calidad}
En esta sección se incluirán las actividades y tareas relacionadas con el aseguramiento de calidad a realizar durante el desarrollo del software. Se incluirán los estándares, prácticas y normas aplicables durante el desarrollo del software.\\

También, deberán recogerse los diferentes tipos de revisiones, verificaciones y validaciones que se van a llevar a cabo, los criterios para la aceptación o rechazo de cada producto y los procedimientos para implementar acciones correctoras o preventivas.


% DESARROLLO
\part{Desarrollo}
\null\vfill
\noindent En esta parte se debe describir el desarrollo del proyecto siguiendo la metodología empleada. Sus capítulos no deben ser una descripción exhaustiva de todos los documentos, diagramas, código fuente y, en general, entregables generados, sino más bien una explicación resumida del desarrollo, estructurada según las etapas principales del proceso de ingeniería. Deben seleccionarse aquellos diagramas, fragmentos de código y secciones de los entregables que sean más significativos para dicha explicación. La totalidad de los entregables resultado del proyecto se ubicarán en los anexos y/o en el material en CD/DVD que acompañe al proyecto.

\chapter{Requisitos del Sistema}
% ------------------------------------------------------------------------------
% Este fichero es parte de la plantilla LaTeX para la realización de Proyectos
% Final de Grado, protegido bajo los términos de la licencia GFDL.
% Para más información, la licencia completa viene incluida en el
% fichero fdl-1.3.tex

% Copyright (C) 2012 SPI-FM. Universidad de Cádiz
% ------------------------------------------------------------------------------

En esta sección se detalla la situación actual de la organización y las necesidades de la misma, que originan el desarrollo o mejora de un sistema informático. Luego se presentan los objetivos y el catálogo de requisitos del nuevo sistema. Finalmente se describen las diferentes alternativas tecnológicas y el análisis de la brecha entre los requisitos planteados y la solución base seleccionada, si aplica.

\section{Situación actual} 
Esta sección debe contener información sobre la situación actual de la organización para la que se va a desarrollar el sistema software.

\subsection{Procesos de Negocio}
Esta sección debe contener información sobre los modelos de procesos de negocio actuales, que suelen ser la base de los modelos de procesos de negocio a implantar.

\subsection{Entorno Tecnológico}
Esta sección debe contener información general sobre el entorno tecnológico en la organización del cliente antes del comienzo del desarrollo del sistema software, incluyendo hardware, redes, software, etc.

\subsection{Fortalezas y Debilidades}
Esta sección debe contener información sobre los aspectos positivos y negativos del negocio actual de la organización para la que se va a desarrollar el sistema software.

\section{Necesidades de Negocio}
Esta sección debe contener información sobre los objetivos de negocio de clientes y usuarios, incluyendo los modelos de procesos de negocio a implantar.

\subsection{Objetivos de Negocio}
Esta sección debe contener los objetivos de negocio que se esperan alcanzar cuando el sistema software a desarrollar esté en producción.

\subsection{Procesos de Negocio}
Esta sección, debe contener los modelos de procesos de negocio a implantar, que normalmente son los modelos de procesos de negocio actuales con ciertas mejoras.

\section{Objetivos del Sistema}
Esta sección debe contener la especificación de los objetivos o requisitos generales del sistema.

\section{Catálogo de Requisitos}
Esta sección debe contener la descripción del conjunto de requisitos específicos del sistema a desarrollar para satisfacer las necesidades de negocio del cliente.

\subsection{Requisitos funcionales}
Descripción completa de la funcionalidad que ofrece el sistema. 

\subsection{Requisitos no funcionales}
Descripción de otros requisitos (relacionados con la calidad del software) que el sistema deberá satisfacer: portabilidad, seguridad, estándares de obligado cumplimiento, accesibilidad, usabilidad, etc.

\subsection{Reglas de negocio}
En el desarrollo del sistema, hay que tener en cuenta las denominadas reglas de negocio, es decir, el conjunto de restricciones, normas o políticas de la organización que deben ser respetadas por el sistema, las cuales suelen ser cambiantes.

\subsection{Requisitos de información}
En esta sección se describen los requisitos de gestión de información (datos) que el sistema debe gestionar.

\section{Alternativas de Solución}
En esta sección, se debe ofrecer un estudio del arte de las diferentes alternativas tecnológicas que permitan satisfacer los requerimientos del sistema, para luego seleccionar (si procede) la herramienta o conjunto de herramientas que utilizaremos como base para el software a desarrollar.

\section{Solución Propuesta}
Si se ha optado por utilizar un software de base, debemos identificar y medir las diferencias entre lo que proporciona este software y los requisitos definidos para el proyecto.\\
El resultado de este análisis permitirá identificar cuáles de éstos requisitos ya están solventados total o parcialmente por el sistema base y cuales tendremos que diseñar e implementar la propuesta de solución.

\chapter{Análisis del Sistema}
% ------------------------------------------------------------------------------
% Este fichero es parte de la plantilla LaTeX para la realización de Proyectos
% Final de Grado, protegido bajo los términos de la licencia GFDL.
% Para más información, la licencia completa viene incluida en el
% fichero fdl-1.3.tex

% Copyright (C) 2012 SPI-FM. Universidad de Cádiz
% ------------------------------------------------------------------------------

Esta sección cubre el análisis del sistema de información a desarrollar, haciendo uso del lenguaje de modelado UML.

\section{Modelo Conceptual}
A partir de los requisitos de información, se desarrollará un diagrama conceptual de clases UML, identificando las clases, atributos, relaciones, restricciones adicionales y reglas de derivación necesarias.

\section{Modelo de Casos de Uso}
A partir de los requisitos funcionales descritos anteriormente, se emplearan los casos de uso como mecanismo para representar las interacciones entre los actores y el sistema bajo estudio. Para cada caso de uso deberá indicarse los actores implicados, las precondiciones y postcondiciones, los pasos que conforman el escenario principal y el conjunto de posibles escenarios alternativos.

\subsection{Actores} 
En este apartado se describirán los diferentes roles que juegan los usuarios que interactúan con el sistema. Los actores pueden ser roles de personas físicas, sistemas externos o incluso el tiempo (eventos temporales).

\section{Modelo de Comportamiento}
A partir de los casos de uso anteriores, se crea el modelo de comportamiento. Para ello, se realizarán los diagramas de secuencia del sistema, donde se identificarán las operaciones o servicios del sistema. Luego, se detallará el contrato de las operaciones identificadas.

\section{Modelo de Interfaz de Usuario}
En esta sección se deberá incluir un prototipo de baja fidelidad o mockup de la interfaz de usuario del sistema. Además, es preciso elaborar un diagrama de navegación, reflejando la secuencia de pantallas a las que tienen acceso los diferentes roles de usuario y la conexión entre éstas.

\chapter{Diseño del Sistema}
% ------------------------------------------------------------------------------
% Este fichero es parte de la plantilla LaTeX para la realización de Proyectos
% Final de Grado, protegido bajo los términos de la licencia GFDL.
% Para más información, la licencia completa viene incluida en el
% fichero fdl-1.3.tex

% Copyright (C) 2012 SPI-FM. Universidad de Cádiz
% ------------------------------------------------------------------------------

En esta sección se recoge la arquitectura general del sistema de información, la parametrización del software base (opcional), el diseño físico de datos, el diseño detallado de componentes software y el diseño detallado de la interfaz de usuario.

\section{Arquitectura del Sistema}
En esta sección se define la arquitectura general del sistema de información, especificando la infraestructura tecnológica necesaria para dar soporte al software y la estructura de los componentes que lo forman.

\subsection{Arquitectura Física}
En este apartado, describimos los principales elementos hardware que forman la arquitectura física de nuestro sistema, recogiendo por un lado los componentes del entorno de producción y los componentes de cliente.\\

Se debe incluir un modelo de despliegue en el cual se describe cómo los elementos software son desplegados en los elementos hardware. También se incluyen las especificaciones y los requisitos del hardware (servidores, etc.), así como de los elementos software (sistemas operativos, servicios, aplicaciones, etc.) necesarios.

\subsection{Arquitectura Lógica}
La arquitectura de diseño especifica la forma en que los artefactos software interactúan entre sí para lograr el comportamiento deseado en el sistema. En esta sección se muestra la comunicación entre el software base seleccionado, los componentes reutilizados y los componentes desarrollados para cumplir los requisitos de la aplicación. También, se recogen los servicios de sistemas externos con los que interactúa nuestro sistema.
Se debe incluir un diagrama de componentes que muestre en un alto nivel de abstracción los artefactos que conforman el sistema.\\

Existen diferentes patrones o estilos arquitectónicos. En los sistemas web de información es común la utilización del patrón Layers (Capas), con el cual estructuramos el sistema en un número apropiado de capas, de forma que todos los componentes de una misma capa trabajan en el mismo nivel de abstracción y los servicios proporcionados por la capa superior utilizan internamente los servicios proporcionados por la capa inmediatamente inferior. Habitualmente se tienen las siguientes capas:

\paragraph*{Capa de presentación (frontend)}
Este grupo de artefactos software conforman la capa de presentación del sistema, incluyendo tanto los componentes de la vista como los elementos de control de la misma.

\paragraph*{Capa de negocio}
Este grupo de artefactos software conforman la capa de negocio del sistema, incluyendo los elementos del modelo de dominio y los servicios (operaciones del sistema).

\paragraph*{Capa de persistencia}
Este grupo de artefactos software conforman la capa de integración del sistema, incluyendo las clases de abstracción para el acceso a datos (BD o sistema de ficheros) o a sistemas heredados.\\

Es común que a la capa de negocio y de datos de los sistemas web, se denomine conjuntamente como backend o modelo de la aplicación.

Opcionalmente, podemos disponer de un conjunto de artefactos software que pueden ser usados por elementos de cualquiera de las capas del sistema y que fundamentalmente proporcionan servicios relacionados con requisitos no funcionales (calidad).

\section{Parametrización del software base}
En esta sección, se detallan las modificaciones a realizar sobre el software base, que son requeridas para la correcta construcción del sistema. En esta sección incluiremos las  actuaciones necesarias sobre la interfaz de administración del sistema, sobre el código fuente o sobre el modelo de datos.

\section{Diseño Físico de Datos}
En esta sección se define la estructura física de datos que utilizará el sistema, a partir del modelo de conceptual de clases, de manera que teniendo presente los requisitos establecidos para el sistema de información y las particularidades del entorno tecnológico, se consiga un acceso eficiente de los datos.
La estructura física se compone de tablas, índices, procedimientos almacenados, secuencias y otros elementos dependientes del SGBD a utilizar.

\section{Diseño detallado de Componentes}
Para cada uno de los módulos funcionales del sistema debemos realizar un diagrama de secuencia, para definir la interacción existente entre las clases de objetos que permitan responder a eventos externos.

\section{Diseño detallado de la Interfaz de Usuario} 
En esta sección se detallarán las interfaces entre el sistema y el usuario, incluyendo un prototipo de alta fidelidad con el diseño de la IU. Se definirá el comportamiento de las diferentes pantallas, indicando qué ocurre en los distintos componentes visuales de la interfaz cuando aparecen y qué acciones se disparan cuando el usuario trabaja con ellas.



\chapter{Construcción del Sistema}
% ------------------------------------------------------------------------------
% Este fichero es parte de la plantilla LaTeX para la realización de Proyectos
% Final de Grado, protegido bajo los términos de la licencia GFDL.
% Para más información, la licencia completa viene incluida en el
% fichero fdl-1.3.tex

% Copyright (C) 2012 SPI-FM. Universidad de Cádiz
% ------------------------------------------------------------------------------

Este capítulo trata sobre todos los aspectos relacionados con la implementación del sistema en código, haciendo uso de un determinado entorno tecnológico.

\section{Entorno de Construcción}
En esta sección se debe indicar el marco tecnológico utilizado para la construcción del sistema: entorno de desarrollo (IDE), lenguaje de programación, herramientas de ayuda a la construcción y despliegue, control de versiones, repositorio de componentes, integración contínua, etc.

\section{Código Fuente}
Organización del código fuente, describiendo la utilidad de los diferentes ficheros y su distribución en paquetes o directorios. Asimismo, se incluirá algún extracto significativo de código fuente que sea de interés para ilustrar algún algoritmo o funcionalidad específica del sistema.

\section{Scripts de Base de datos}
Organización del código fuente, describiendo la utilidad de los diferentes ficheros y su distribución en paquetes o directorios. Asimismo, se incluirá el script de algún disparador o un procedimiento almacenado, que sea de interés para ilustrar algún aspecto concreto de la gestión de la base de datos.


\chapter{Pruebas del Sistema}
% ------------------------------------------------------------------------------
% Este fichero es parte de la plantilla LaTeX para la realización de Proyectos
% Final de Grado, protegido bajo los términos de la licencia GFDL.
% Para más información, la licencia completa viene incluida en el
% fichero fdl-1.3.tex

% Copyright (C) 2012 SPI-FM. Universidad de Cádiz
% ------------------------------------------------------------------------------

En este capítulo se presenta el plan de pruebas del sistema de información, incluyendo los diferentes tipos de pruebas que se han llevado a cabo, ya sean manuales (mediante listas de comprobación) o automatizadas mediante algún software específico de pruebas.

\section{Estrategia}
En esta sección se debe incluir el alcance de las pruebas, hasta donde se pretende llegar con ellas, si se registrarán todas o sólo aquellas de un cierto tipo y cómo se interpretarán y evaluarán los resultados.
También, se incluirá el procedimiento a seguir para las pruebas de regresión, esto es, la repetición de ciertas pruebas para comprobar que nuevos cambios que se vayan introduciendo no originen errores en el software ya probado.

\section{Entorno de Pruebas}
Incluir en este apartado los requisitos de los entornos hardware/software donde se ejecutarán las pruebas.

\section{Roles}
Describir en esa sección cuáles serán los perfiles y participantes necesarios para la ejecución de cada uno de los niveles de prueba.

\section{Niveles de Pruebas}
En este sección se documentan los diferentes tipos de pruebas que se han llevado a cabo, ya sean manuales o automatizadas mediante algún software específico de pruebas.

\subsection{Pruebas Unitarias}
Las pruebas unitarias tienen por objetivo localizar errores en cada nuevo artefacto software desarrollado, antes que se produzca la integración con el resto de artefactos del sistema.

\subsection{Pruebas de Integración}
Este tipo de pruebas tienen por objetivo localizar errores en módulos o subsistemas completos, analizando la interacción entre varios artefactos software.

\subsection{Pruebas de Sistema}
En esta actividad se realizan las pruebas de sistema de modo que se asegure que el sistema cumple con todos los requisitos establecidos: funcionales, de almacenamiento, reglas de negocio y no funcionales. Se suelen desarrollar en un entorno específico para pruebas.

\subsubsection{Pruebas Funcionales}
Con estas pruebas se analiza el buen funcionamiento de la implementación de los flujos normales y alternativos de los distintos casos de uso del sistema.

\subsubsection{Pruebas No Funcionales}
Estas pruebas pretenden comprobar el funcionamiento del sistema, con respecto a los requisitos no funcionales identificados: eficiencia, seguridad, etc.

\subsection{Pruebas de Aceptación}
El objetivo de estas pruebas es demostrar que el producto está listo para el paso a producción. Suelen ser las mismas pruebas que se realizaron anteriormente pero en el entorno de producción. En estas pruebas, es importante la participación del cliente final.

% EPILOGO
\part{Epílogo}
\null\vfill
\noindent En esta última parte quedarán recogidas las conclusiones y los manuales necesarios para el manejo de la aplicación resultado del desarrollo. Si se ha realizado algún tipo de evaluación de la solución proporcionada, más allá de las pruebas del sistema, también deberá venir recogida en un capítulo separado dentro de esta parte. Pueden consultarse diversos tipos de evaluaciones sobre sistemas de información en \cite{hevner2004}: casos de estudio, análisis estático, análisis dinámico, simulación, experimento controlado, etc.
\vfill

\chapter{Manual de implantación y explotación}
% ------------------------------------------------------------------------------
% Este fichero es parte de la plantilla LaTeX para la realización de Proyectos
% Final de Grado, protegido bajo los términos de la licencia GFDL.
% Para más información, la licencia completa viene incluida en el
% fichero fdl-1.3.tex

% Copyright (C) 2012 SPI-FM. Universidad de Cádiz
% ------------------------------------------------------------------------------

Las instrucciones de instalación y explotación del sistema se detallan a continuación.

\section{Introducción}
Resumen de los principales objetivos, ámbito y alcance del software desarrollado.

\section{Requisitos previos}
Requisitos hardware y software para la correcta instalación del sistema.

\section{Inventario de componentes}
Lista de los componentes hardware y software que se incluyen en la versión del producto.

\section{Procedimientos de instalación}
Procedimientos de instalación y configuración de cada componente hardware y software (base y desarrollado) para asegurar la correcta instalación y explotación del sistema, así como aquellos procedimientos necesarios de migración/carga de datos.

\section{Pruebas de implantación}
Descripción de las pruebas a realizar después de la instalación del sistema. 

\section{Procedimientos de operación y nivel de servicio}
Procedimientos necesarios para asegurar el correcto funcionamiento, rendimiento, disponibilidad y seguridad del sistema: back-ups, chequeo de logs, etc. También, es preciso indicar claramente aquellas actuaciones precisas necesarias para el mantenimiento preventivo del sistema y así prevenir posibles fallos en el mismo. 


\chapter{Manual de usuario}
% ------------------------------------------------------------------------------
% Este fichero es parte de la plantilla LaTeX para la realización de Proyectos
% Final de Grado, protegido bajo los términos de la licencia GFDL.
% Para más información, la licencia completa viene incluida en el
% fichero fdl-1.3.tex

% Copyright (C) 2012 SPI-FM. Universidad de Cádiz
% ------------------------------------------------------------------------------

Las instrucciones de uso del sistema se detallan a continuación.

\section{Introducción}
Resumen de los principales objetivos, ámbito y alcance del software desarrollado.

\section{Características}
Recopilación de las principales funcionalidades del sistema.

\section{Requisitos previos}
Requisitos hardware y software para el correcto uso del sistema.

\section{Uso del sistema}
Describir todos los aspectos necesarios para una utilización efectiva y eficiente del sistema por parte de los usuarios.





\chapter{Conclusiones}
% ------------------------------------------------------------------------------
% Este fichero es parte de la plantilla LaTeX para la realización de Proyectos
% Final de Grado, protegido bajo los términos de la licencia GFDL.
% Para más información, la licencia completa viene incluida en el
% fichero fdl-1.3.tex

% Copyright (C) 2012 SPI-FM. Universidad de Cádiz
% ------------------------------------------------------------------------------

En este último capítulo se detallan las lecciones aprendidas tras el desarrollo del presente proyecto y se identifican las posibles oportunidades de mejora sobre el software desarrollado.

\section{Objetivos alcanzados}
Este apartado debe resumir los objetivos generales y específicos alcanzados, relacionándolos con todo lo descrito en el capítulo de introducción.\\

\section{Lecciones aprendidas}
A continuación, se detallan las buenas prácticas adquiridas, tanto tecnológicas como procedimentales, así como cualquier otro aspecto de interés.\\
Resumir cuantitativamente el tiempo y esfuerzo dedicados al proyecto a lo largo de su desarrollo que escribir un sencillo 'he trabajado mucho en este proyecto'.

\section{Trabajo futuro}
En esta sección, se presentan las diversas áreas u oportunidades de mejora detectadas durante el desarrollo del proyecto y que podrán ser abarcadas en futuras versiones del software.\\

Los elementos aquí descritos deben estar en relación con lo relatado en el apartado de objetivos y alcance del proyecto descritos en la introducción.



\chapter*{\bibname}
\addcontentsline{toc}{chapter}{\bibname}
%\renewcommand{\bibname}{}

%\input{./bibliografia}

\begingroup
  \def\chapter*#1{}
\renewcommand{\bibname}{}
% Bibliografía con BibTeX
\bibliographystyle{apalike}
\bibliography{bibliografia}

\backmatter

% ------------------------------------------------------------------------------
% Este fichero es parte de la plantilla LaTeX para la realización de Proyectos
% Final de Grado, protegido bajo los términos de la licencia GFDL.
% Para más información, la licencia completa viene incluida en el
% fichero fdl-1.3.tex

% Copyright (C) 2012 SPI-FM. Universidad de Cádiz
% ------------------------------------------------------------------------------


\chapter*{\rlap{Información sobre Licencia}}
\phantomsection  % so hyperref creates bookmarks
\addcontentsline{toc}{chapter}{Información sobre Licencia}
%\label{label_fdl}

 \begin{center}

       Información sobre Licencia


\end{center}

Incluir aquí la información relativa a la licencia seleccionada para la documentación y software del presente proyecto.

%GNU % ------------------------------------------------------------------------------
% Este fichero es parte de la plantilla LaTeX para la realización de Proyectos
% Final de Grado, protegido bajo los términos de la licencia GFDL.
% Para más información, la licencia completa viene incluida en el
% fichero fdl-1.3.tex

% Copyright (C) 2012 SPI-FM. Universidad de Cádiz
% ------------------------------------------------------------------------------


\chapter*{\rlap{GNU Free Documentation License}}
\phantomsection  % so hyperref creates bookmarks
\addcontentsline{toc}{chapter}{GNU Free Documentation License}
%\label{label_fdl}

 \begin{center}

       Version 1.3, 3 November 2008


 Copyright \copyright{} 2000, 2001, 2002, 2007, 2008  Free Software Foundation, Inc.
 
 \bigskip
 
     <http://fsf.org/>
  
 \bigskip
 
 Everyone is permitted to copy and distribute verbatim copies
 of this license document, but changing it is not allowed.
\end{center}


\begin{center}
{\bf\large Preamble}
\end{center}

The purpose of this License is to make a manual, textbook, or other
functional and useful document ``free'' in the sense of freedom: to
assure everyone the effective freedom to copy and redistribute it,
with or without modifying it, either commercially or noncommercially.
Secondarily, this License preserves for the author and publisher a way
to get credit for their work, while not being considered responsible
for modifications made by others.

This License is a kind of ``copyleft'', which means that derivative
works of the document must themselves be free in the same sense.  It
complements the GNU General Public License, which is a copyleft
license designed for free software.

We have designed this License in order to use it for manuals for free
software, because free software needs free documentation: a free
program should come with manuals providing the same freedoms that the
software does.  But this License is not limited to software manuals;
it can be used for any textual work, regardless of subject matter or
whether it is published as a printed book.  We recommend this License
principally for works whose purpose is instruction or reference.


\begin{center}
{\Large\bf 1. APPLICABILITY AND DEFINITIONS\par}
\phantomsection
\addcontentsline{toc}{section}{1. APPLICABILITY AND DEFINITIONS}
\end{center}

This License applies to any manual or other work, in any medium, that
contains a notice placed by the copyright holder saying it can be
distributed under the terms of this License.  Such a notice grants a
world-wide, royalty-free license, unlimited in duration, to use that
work under the conditions stated herein.  The ``\textbf{Document}'', below,
refers to any such manual or work.  Any member of the public is a
licensee, and is addressed as ``\textbf{you}''.  You accept the license if you
copy, modify or distribute the work in a way requiring permission
under copyright law.

A ``\textbf{Modified Version}'' of the Document means any work containing the
Document or a portion of it, either copied verbatim, or with
modifications and/or translated into another language.

A ``\textbf{Secondary Section}'' is a named appendix or a front-matter section of
the Document that deals exclusively with the relationship of the
publishers or authors of the Document to the Document's overall subject
(or to related matters) and contains nothing that could fall directly
within that overall subject.  (Thus, if the Document is in part a
textbook of mathematics, a Secondary Section may not explain any
mathematics.)  The relationship could be a matter of historical
connection with the subject or with related matters, or of legal,
commercial, philosophical, ethical or political position regarding
them.

The ``\textbf{Invariant Sections}'' are certain Secondary Sections whose titles
are designated, as being those of Invariant Sections, in the notice
that says that the Document is released under this License.  If a
section does not fit the above definition of Secondary then it is not
allowed to be designated as Invariant.  The Document may contain zero
Invariant Sections.  If the Document does not identify any Invariant
Sections then there are none.

The ``\textbf{Cover Texts}'' are certain short passages of text that are listed,
as Front-Cover Texts or Back-Cover Texts, in the notice that says that
the Document is released under this License.  A Front-Cover Text may
be at most 5 words, and a Back-Cover Text may be at most 25 words.

A ``\textbf{Transparent}'' copy of the Document means a machine-readable copy,
represented in a format whose specification is available to the
general public, that is suitable for revising the document
straightforwardly with generic text editors or (for images composed of
pixels) generic paint programs or (for drawings) some widely available
drawing editor, and that is suitable for input to text formatters or
for automatic translation to a variety of formats suitable for input
to text formatters.  A copy made in an otherwise Transparent file
format whose markup, or absence of markup, has been arranged to thwart
or discourage subsequent modification by readers is not Transparent.
An image format is not Transparent if used for any substantial amount
of text.  A copy that is not ``Transparent'' is called ``\textbf{Opaque}''.

Examples of suitable formats for Transparent copies include plain
ASCII without markup, Texinfo input format, LaTeX input format, SGML
or XML using a publicly available DTD, and standard-conforming simple
HTML, PostScript or PDF designed for human modification.  Examples of
transparent image formats include PNG, XCF and JPG.  Opaque formats
include proprietary formats that can be read and edited only by
proprietary word processors, SGML or XML for which the DTD and/or
processing tools are not generally available, and the
machine-generated HTML, PostScript or PDF produced by some word
processors for output purposes only.

The ``\textbf{Title Page}'' means, for a printed book, the title page itself,
plus such following pages as are needed to hold, legibly, the material
this License requires to appear in the title page.  For works in
formats which do not have any title page as such, ``Title Page'' means
the text near the most prominent appearance of the work's title,
preceding the beginning of the body of the text.

The ``\textbf{publisher}'' means any person or entity that distributes
copies of the Document to the public.

A section ``\textbf{Entitled XYZ}'' means a named subunit of the Document whose
title either is precisely XYZ or contains XYZ in parentheses following
text that translates XYZ in another language.  (Here XYZ stands for a
specific section name mentioned below, such as ``\textbf{Acknowledgements}'',
``\textbf{Dedications}'', ``\textbf{Endorsements}'', or ``\textbf{History}''.)  
To ``\textbf{Preserve the Title}''
of such a section when you modify the Document means that it remains a
section ``Entitled XYZ'' according to this definition.

The Document may include Warranty Disclaimers next to the notice which
states that this License applies to the Document.  These Warranty
Disclaimers are considered to be included by reference in this
License, but only as regards disclaiming warranties: any other
implication that these Warranty Disclaimers may have is void and has
no effect on the meaning of this License.


\begin{center}
{\Large\bf 2. VERBATIM COPYING\par}
\phantomsection
\addcontentsline{toc}{section}{2. VERBATIM COPYING}
\end{center}

You may copy and distribute the Document in any medium, either
commercially or noncommercially, provided that this License, the
copyright notices, and the license notice saying this License applies
to the Document are reproduced in all copies, and that you add no other
conditions whatsoever to those of this License.  You may not use
technical measures to obstruct or control the reading or further
copying of the copies you make or distribute.  However, you may accept
compensation in exchange for copies.  If you distribute a large enough
number of copies you must also follow the conditions in section~3.

You may also lend copies, under the same conditions stated above, and
you may publicly display copies.


\begin{center}
{\Large\bf 3. COPYING IN QUANTITY\par}
\phantomsection
\addcontentsline{toc}{section}{3. COPYING IN QUANTITY}
\end{center}


If you publish printed copies (or copies in media that commonly have
printed covers) of the Document, numbering more than 100, and the
Document's license notice requires Cover Texts, you must enclose the
copies in covers that carry, clearly and legibly, all these Cover
Texts: Front-Cover Texts on the front cover, and Back-Cover Texts on
the back cover.  Both covers must also clearly and legibly identify
you as the publisher of these copies.  The front cover must present
the full title with all words of the title equally prominent and
visible.  You may add other material on the covers in addition.
Copying with changes limited to the covers, as long as they preserve
the title of the Document and satisfy these conditions, can be treated
as verbatim copying in other respects.

If the required texts for either cover are too voluminous to fit
legibly, you should put the first ones listed (as many as fit
reasonably) on the actual cover, and continue the rest onto adjacent
pages.

If you publish or distribute Opaque copies of the Document numbering
more than 100, you must either include a machine-readable Transparent
copy along with each Opaque copy, or state in or with each Opaque copy
a computer-network location from which the general network-using
public has access to download using public-standard network protocols
a complete Transparent copy of the Document, free of added material.
If you use the latter option, you must take reasonably prudent steps,
when you begin distribution of Opaque copies in quantity, to ensure
that this Transparent copy will remain thus accessible at the stated
location until at least one year after the last time you distribute an
Opaque copy (directly or through your agents or retailers) of that
edition to the public.

It is requested, but not required, that you contact the authors of the
Document well before redistributing any large number of copies, to give
them a chance to provide you with an updated version of the Document.


\begin{center}
{\Large\bf 4. MODIFICATIONS\par}
\phantomsection
\addcontentsline{toc}{section}{4. MODIFICATIONS}
\end{center}

You may copy and distribute a Modified Version of the Document under
the conditions of sections 2 and 3 above, provided that you release
the Modified Version under precisely this License, with the Modified
Version filling the role of the Document, thus licensing distribution
and modification of the Modified Version to whoever possesses a copy
of it.  In addition, you must do these things in the Modified Version:

\begin{itemize}
\item[A.] 
   Use in the Title Page (and on the covers, if any) a title distinct
   from that of the Document, and from those of previous versions
   (which should, if there were any, be listed in the History section
   of the Document).  You may use the same title as a previous version
   if the original publisher of that version gives permission.
   
\item[B.]
   List on the Title Page, as authors, one or more persons or entities
   responsible for authorship of the modifications in the Modified
   Version, together with at least five of the principal authors of the
   Document (all of its principal authors, if it has fewer than five),
   unless they release you from this requirement.
   
\item[C.]
   State on the Title page the name of the publisher of the
   Modified Version, as the publisher.
   
\item[D.]
   Preserve all the copyright notices of the Document.
   
\item[E.]
   Add an appropriate copyright notice for your modifications
   adjacent to the other copyright notices.
   
\item[F.]
   Include, immediately after the copyright notices, a license notice
   giving the public permission to use the Modified Version under the
   terms of this License, in the form shown in the Addendum below.
   
\item[G.]
   Preserve in that license notice the full lists of Invariant Sections
   and required Cover Texts given in the Document's license notice.
   
\item[H.]
   Include an unaltered copy of this License.
   
\item[I.]
   Preserve the section Entitled ``History'', Preserve its Title, and add
   to it an item stating at least the title, year, new authors, and
   publisher of the Modified Version as given on the Title Page.  If
   there is no section Entitled ``History'' in the Document, create one
   stating the title, year, authors, and publisher of the Document as
   given on its Title Page, then add an item describing the Modified
   Version as stated in the previous sentence.
   
\item[J.]
   Preserve the network location, if any, given in the Document for
   public access to a Transparent copy of the Document, and likewise
   the network locations given in the Document for previous versions
   it was based on.  These may be placed in the ``History'' section.
   You may omit a network location for a work that was published at
   least four years before the Document itself, or if the original
   publisher of the version it refers to gives permission.
   
\item[K.]
   For any section Entitled ``Acknowledgements'' or ``Dedications'',
   Preserve the Title of the section, and preserve in the section all
   the substance and tone of each of the contributor acknowledgements
   and/or dedications given therein.
   
\item[L.]
   Preserve all the Invariant Sections of the Document,
   unaltered in their text and in their titles.  Section numbers
   or the equivalent are not considered part of the section titles.
   
\item[M.]
   Delete any section Entitled ``Endorsements''.  Such a section
   may not be included in the Modified Version.
   
\item[N.]
   Do not retitle any existing section to be Entitled ``Endorsements''
   or to conflict in title with any Invariant Section.
   
\item[O.]
   Preserve any Warranty Disclaimers.
\end{itemize}

If the Modified Version includes new front-matter sections or
appendices that qualify as Secondary Sections and contain no material
copied from the Document, you may at your option designate some or all
of these sections as invariant.  To do this, add their titles to the
list of Invariant Sections in the Modified Version's license notice.
These titles must be distinct from any other section titles.

You may add a section Entitled ``Endorsements'', provided it contains
nothing but endorsements of your Modified Version by various
parties---for example, statements of peer review or that the text has
been approved by an organization as the authoritative definition of a
standard.

You may add a passage of up to five words as a Front-Cover Text, and a
passage of up to 25 words as a Back-Cover Text, to the end of the list
of Cover Texts in the Modified Version.  Only one passage of
Front-Cover Text and one of Back-Cover Text may be added by (or
through arrangements made by) any one entity.  If the Document already
includes a cover text for the same cover, previously added by you or
by arrangement made by the same entity you are acting on behalf of,
you may not add another; but you may replace the old one, on explicit
permission from the previous publisher that added the old one.

The author(s) and publisher(s) of the Document do not by this License
give permission to use their names for publicity for or to assert or
imply endorsement of any Modified Version.


\begin{center}
{\Large\bf 5. COMBINING DOCUMENTS\par}
\phantomsection
\addcontentsline{toc}{section}{5. COMBINING DOCUMENTS}
\end{center}


You may combine the Document with other documents released under this
License, under the terms defined in section~4 above for modified
versions, provided that you include in the combination all of the
Invariant Sections of all of the original documents, unmodified, and
list them all as Invariant Sections of your combined work in its
license notice, and that you preserve all their Warranty Disclaimers.

The combined work need only contain one copy of this License, and
multiple identical Invariant Sections may be replaced with a single
copy.  If there are multiple Invariant Sections with the same name but
different contents, make the title of each such section unique by
adding at the end of it, in parentheses, the name of the original
author or publisher of that section if known, or else a unique number.
Make the same adjustment to the section titles in the list of
Invariant Sections in the license notice of the combined work.

In the combination, you must combine any sections Entitled ``History''
in the various original documents, forming one section Entitled
``History''; likewise combine any sections Entitled ``Acknowledgements'',
and any sections Entitled ``Dedications''.  You must delete all sections
Entitled ``Endorsements''.

\begin{center}
{\Large\bf 6. COLLECTIONS OF DOCUMENTS\par}
\phantomsection
\addcontentsline{toc}{section}{6. COLLECTIONS OF DOCUMENTS}
\end{center}

You may make a collection consisting of the Document and other documents
released under this License, and replace the individual copies of this
License in the various documents with a single copy that is included in
the collection, provided that you follow the rules of this License for
verbatim copying of each of the documents in all other respects.

You may extract a single document from such a collection, and distribute
it individually under this License, provided you insert a copy of this
License into the extracted document, and follow this License in all
other respects regarding verbatim copying of that document.


\begin{center}
{\Large\bf 7. AGGREGATION WITH INDEPENDENT WORKS\par}
\phantomsection
\addcontentsline{toc}{section}{7. AGGREGATION WITH INDEPENDENT WORKS}
\end{center}


A compilation of the Document or its derivatives with other separate
and independent documents or works, in or on a volume of a storage or
distribution medium, is called an ``aggregate'' if the copyright
resulting from the compilation is not used to limit the legal rights
of the compilation's users beyond what the individual works permit.
When the Document is included in an aggregate, this License does not
apply to the other works in the aggregate which are not themselves
derivative works of the Document.

If the Cover Text requirement of section~3 is applicable to these
copies of the Document, then if the Document is less than one half of
the entire aggregate, the Document's Cover Texts may be placed on
covers that bracket the Document within the aggregate, or the
electronic equivalent of covers if the Document is in electronic form.
Otherwise they must appear on printed covers that bracket the whole
aggregate.


\begin{center}
{\Large\bf 8. TRANSLATION\par}
\phantomsection
\addcontentsline{toc}{section}{8. TRANSLATION}
\end{center}


Translation is considered a kind of modification, so you may
distribute translations of the Document under the terms of section~4.
Replacing Invariant Sections with translations requires special
permission from their copyright holders, but you may include
translations of some or all Invariant Sections in addition to the
original versions of these Invariant Sections.  You may include a
translation of this License, and all the license notices in the
Document, and any Warranty Disclaimers, provided that you also include
the original English version of this License and the original versions
of those notices and disclaimers.  In case of a disagreement between
the translation and the original version of this License or a notice
or disclaimer, the original version will prevail.

If a section in the Document is Entitled ``Acknowledgements'',
``Dedications'', or ``History'', the requirement (section~4) to Preserve
its Title (section~1) will typically require changing the actual
title.


\begin{center}
{\Large\bf 9. TERMINATION\par}
\phantomsection
\addcontentsline{toc}{section}{9. TERMINATION}
\end{center}


You may not copy, modify, sublicense, or distribute the Document
except as expressly provided under this License.  Any attempt
otherwise to copy, modify, sublicense, or distribute it is void, and
will automatically terminate your rights under this License.

However, if you cease all violation of this License, then your license
from a particular copyright holder is reinstated (a) provisionally,
unless and until the copyright holder explicitly and finally
terminates your license, and (b) permanently, if the copyright holder
fails to notify you of the violation by some reasonable means prior to
60 days after the cessation.

Moreover, your license from a particular copyright holder is
reinstated permanently if the copyright holder notifies you of the
violation by some reasonable means, this is the first time you have
received notice of violation of this License (for any work) from that
copyright holder, and you cure the violation prior to 30 days after
your receipt of the notice.

Termination of your rights under this section does not terminate the
licenses of parties who have received copies or rights from you under
this License.  If your rights have been terminated and not permanently
reinstated, receipt of a copy of some or all of the same material does
not give you any rights to use it.


\begin{center}
{\Large\bf 10. FUTURE REVISIONS OF THIS LICENSE\par}
\phantomsection
\addcontentsline{toc}{section}{10. FUTURE REVISIONS OF THIS LICENSE}
\end{center}


The Free Software Foundation may publish new, revised versions
of the GNU Free Documentation License from time to time.  Such new
versions will be similar in spirit to the present version, but may
differ in detail to address new problems or concerns.  See
http://www.gnu.org/copyleft/.

Each version of the License is given a distinguishing version number.
If the Document specifies that a particular numbered version of this
License ``or any later version'' applies to it, you have the option of
following the terms and conditions either of that specified version or
of any later version that has been published (not as a draft) by the
Free Software Foundation.  If the Document does not specify a version
number of this License, you may choose any version ever published (not
as a draft) by the Free Software Foundation.  If the Document
specifies that a proxy can decide which future versions of this
License can be used, that proxy's public statement of acceptance of a
version permanently authorizes you to choose that version for the
Document.


\begin{center}
{\Large\bf 11. RELICENSING\par}
\phantomsection
\addcontentsline{toc}{section}{11. RELICENSING}
\end{center}


``Massive Multiauthor Collaboration Site'' (or ``MMC Site'') means any
World Wide Web server that publishes copyrightable works and also
provides prominent facilities for anybody to edit those works.  A
public wiki that anybody can edit is an example of such a server.  A
``Massive Multiauthor Collaboration'' (or ``MMC'') contained in the
site means any set of copyrightable works thus published on the MMC
site.

``CC-BY-SA'' means the Creative Commons Attribution-Share Alike 3.0
license published by Creative Commons Corporation, a not-for-profit
corporation with a principal place of business in San Francisco,
California, as well as future copyleft versions of that license
published by that same organization.

``Incorporate'' means to publish or republish a Document, in whole or
in part, as part of another Document.

An MMC is ``eligible for relicensing'' if it is licensed under this
License, and if all works that were first published under this License
somewhere other than this MMC, and subsequently incorporated in whole
or in part into the MMC, (1) had no cover texts or invariant sections,
and (2) were thus incorporated prior to November 1, 2008.

The operator of an MMC Site may republish an MMC contained in the site
under CC-BY-SA on the same site at any time before August 1, 2009,
provided the MMC is eligible for relicensing.


\begin{center}
{\Large\bf ADDENDUM: How to use this License for your documents\par}
\phantomsection
\addcontentsline{toc}{section}{ADDENDUM: How to use this License for your documents}
\end{center}

To use this License in a document you have written, include a copy of
the License in the document and put the following copyright and
license notices just after the title page:

\bigskip
\begin{quote}
    Copyright \copyright{}  YEAR  YOUR NAME.
    Permission is granted to copy, distribute and/or modify this document
    under the terms of the GNU Free Documentation License, Version 1.3
    or any later version published by the Free Software Foundation;
    with no Invariant Sections, no Front-Cover Texts, and no Back-Cover Texts.
    A copy of the license is included in the section entitled ``GNU
    Free Documentation License''.
\end{quote}
\bigskip
    
If you have Invariant Sections, Front-Cover Texts and Back-Cover Texts,
replace the ``with \dots\ Texts.'' line with this:

\bigskip
\begin{quote}
    with the Invariant Sections being LIST THEIR TITLES, with the
    Front-Cover Texts being LIST, and with the Back-Cover Texts being LIST.
\end{quote}
\bigskip
    
If you have Invariant Sections without Cover Texts, or some other
combination of the three, merge those two alternatives to suit the
situation.

If your document contains nontrivial examples of program code, we
recommend releasing these examples in parallel under your choice of
free software license, such as the GNU General Public License,
to permit their use in free software.

%---------------------------------------------------------------------


\end{document}
